One deduces that the map $\mu\circ(\phi\circ\mathrm{pr}_1,\phi\circ\mathrm{pr}_2):G\times G\to\mathbf W(U)$ corresponds to the composite map $\beta$ below:
\[
\xymatrix@C=30pt@R=18pt{
&H\otimes H\ar[r]^\beta\ar[dl]_{\mathrm{id}_H\circ\sigma_{23}\circ\mathrm{id}_H}&U\\
H\otimes H\otimes H\otimes H\ar[dr]_{(\mathrm{id}_H\otimes\varepsilon_H)\otimes(\varepsilon_H\otimes\mathrm{id}_H)}&&U\otimes U\ar[u]_{m_U}\\
&H\otimes H\ar[ur]_{\phi\otimes\phi}&.}
\]
% typo?: phi instead of theta in the composed functor maps, and id_H circle sigma_23 circle id_H on the first arrow, retained as printed.
Finally, $\theta$ is compatible with the laws of $G$ and $\mathbf W(U)^\times$ if and only if $\theta\circ m_G$ equals $\mu\circ(\phi\circ\mathrm{pr}_1,\phi\circ\mathrm{pr}_2)$, which is equivalent, by the foregoing, to $\phi\circ m_H=\beta$. Now it is clear that $(\phi\circ m_H)(x\otimes y)=\phi(xy)$, and one easily sees that $\beta(x\otimes y)=\phi(x)\phi(y)$.

\numpara{2.4}Let us now return to any pseudocompact ring $k$ to apply to formal groups the results of 1.4--1.5 on passage to the quotient by a topologically flat equivalence relation.
\nde{We have detailed the original in what follows.}Let $u:H\to G$ be a monomorphism of formal $k$-groups, $\mu:G\times G\to G$ the ``multiplication'' morphism of $G$, and $\lambda$ the composite morphism
\[
\lambda:G\times H\xrightarrow{\mathrm{id}_G\times u}G\times G\xrightarrow{\mu}G.
\]
Since $u$ is a monomorphism, the pair
\[
\xymatrix{G\times H\ar@<.5ex>[r]^{\mathrm{pr}_1}\ar@<-.5ex>[r]_\lambda&G}
\]
is an equivalence pair in $\mathbf{Vaf}/k$ (cf. V, 2.b)). Recall (cf. 1.2.C) that the cokernel $G/H$ of this pair is defined as follows.
Let $\OO(G)$ and $\OO(H)$ be the affine algebras of $G$ and $H$, $\Delta:\OO(G)\to\OO(G)\widehat\otimes_k\OO(G)$ the diagonal morphism of $\OO(G)$, and $I$ the kernel of the morphism $u^\natural:\OO(G)\to\OO(H)$. (By Proposition 1.3 one knows that $u^\natural$ induces an isomorphism $\OO(G)/I\isomto\OO(H)$.) Then the affine algebra $\OO(G/H)$ of $G/H$ is the kernel of the pair of morphisms
\[
\xymatrix@C=54pt{\OO(G)\ar@<.5ex>[r]^{\tau_1}\ar@<-.5ex>[r]_{(\mathrm{id}\widehat\otimes u^\natural)\Delta}&\OO(G)\widehat\otimes_k\OO(H),}
\]
where $\tau_1(x)=x\widehat\otimes1$, \ie
\[
\OO(G/H)=\{x\in\OO(G)\mid\Delta(x)-x\widehat\otimes1\in\OO(G)\widehat\otimes I\}.
\]
If, moreover, $H$ is \emph{topologically flat over $k$}, $\mathrm{pr}_1$ is topologically flat and from Theorem 1.4 one deduces the following theorem.
\begin{theorem}{}
Let $u:H\to G$ be a monomorphism of formal $k$-groups. Suppose that $H$ is topologically flat over $k$. Then the projection $p:G\to G/H$ is surjective and topologically flat, one has an isomorphism
\begin{equation}\tag{*}
G\times H\isomto G\times_{G/H}G,
\end{equation}
and $G/H$ represents the quotient sheaf for the flat topology.
Consequently, $G/H$ is equipped with a canonical structure of an object with group of operators $G$, such that $p:G\to G/H$ is a morphism of objects with operators. If, moreover, $u$ identifies $H$ with a normal subgroup of $G$, then $G/H$ is equipped with a canonical formal $k$-group structure such that $p:G\to G/H$ is a morphism of formal $k$-groups, and $H$ is the kernel of $p$.
\end{theorem}
Indeed, the first assertion follows from 1.4; the other two from IV, Corollaries 5.2.2 and 5.2.4.
\begin{corollary}{}\nde{We have made explicit this corollary, which is used, for example, in 5.2.1/5.2.3.}
Let $G$ be a formal $k$-group, $H$ a formal subgroup of $G$, $A$ (resp. $A/J,B$) the affine algebra of $G$ (resp. $H,G/H$), $I_A$ the augmentation ideal of $A$, and $I_B=B\cap I_A$. Suppose that $H$ is topologically flat over $k$. Then $J$ equals $\overline{AI_B}$, the closed ideal generated by $I_B$.
\end{corollary}
Indeed, the projection $B\to B/I_B$ corresponds to the ``identity section'' $e:\Spf(k)\to G/H$ of $G/H$. By (*), $H$ is identified with the fibered product $\Spf(k)\times_{G/H}G$, so its affine algebra $A/J$ is identified with $(B/I_B)\widehat\otimes_B A\simeq A/\overline{AI_B}$.

\numpara{2.4.A}\nde{We have added paragraphs 2.4.A and 2.4.B, which will be useful in 5.1.3 and 2.9 respectively.}Let $G,Q$ be topologically flat formal $k$-groups; suppose that there exist homomorphisms $\sigma:Q\to G$ and $\pi:G\to Q$ such that $\pi\circ\sigma=\mathrm{id}_Q$. In particular, $\sigma$ is a monomorphism, so $Q$ is a formal subgroup of $G$ (cf. Remark 1.3). Let $N=\Ker(\pi)$ and $\sigma'$ be the inclusion $N\hookrightarrow G$. Then $G$ is the \emph{semidirect product} of $N$ by $Q$ (cf. I, 2.3.8), \ie for every $B\in\Ob\mathbf{Alf}/k$, identifying $N(B)$ and $Q(B)$ with their images in $G(B)$ under $\sigma'$ and $\sigma$, $N(B)$ is a normal subgroup of $G(B)$ and the map
\begin{equation}\tag{1}
\mu:N(B)\times Q(B)\to G(B),\qquad(x,q)\mto xq
\end{equation}
is bijective. Then the morphism of formal $k$-varieties
\begin{equation}\tag{2}
\theta:G(B)\to N(B),\qquad g\mto g\cdot\sigma\pi(g^{-1})
\end{equation}
is a retraction of $\sigma'$, the inverse of $\mu$ is the map
\begin{equation}\tag{3}
g\mto(\theta(g),\pi(g)),
\end{equation}
and $\theta\circ\sigma':N\to G/Q$ is an isomorphism of formal $k$-varieties. In particular, $N$ is topologically flat over $k$, by 1.4(ii).
% typo?: theta circle sigma' is printed with target G/Q; retained.
Denote by $\alpha$ (resp. $\beta$) the map from $Q\times N$ (resp. $G\times N=N\times Q\times N$) to $N$ defined set-theoretically by $\alpha(q,y)=qyq^{-1}$ (resp. $\beta(x,q,y)=x\alpha(q,y)$). Then one has the commutative diagram below:
% typo: missing closing parenthesis after beta's definition -> supplied.
\begin{equation}\tag{4}
\xymatrix@C=28pt{
G\times G\ar[r]^{\mathrm{id}\times\theta}\ar[d]_{m_G}&G\times N\ar[r]^{\theta\times\mathrm{id}}\ar[d]^\beta&N\times N\ar[dl]^{m_N}\\
G\ar[r]_\theta&N&.}
\end{equation}
This can be expressed as follows in terms of the affine algebras $A,A_0,A'$ of $G,Q,N$ (cf. 5.1.3 below). Let $\rho':A\to A'$, $\rho:A\to A_0$, and $\tau:A_0\to A$ be the homomorphisms of $k$-bialgebras corresponding to $\sigma',\sigma,\pi$, and let $I=\Ker(\rho)$. Then, by the preceding corollary, $A'$ is identified with $A/\overline{A\tau(J_0)}$, where $J_0$ denotes the augmentation ideal of $A_0$.
On the other hand, let $B$ be the affine algebra of $G/Q$; it is the kernel of the pair of morphisms
\[
\xymatrix@C=58pt{A\ar@<.5ex>[r]^{\tau_1}\ar@<-.5ex>[r]_{(\mathrm{id}\widehat\otimes\rho)\Delta_A}&A\widehat\otimes_k A_0,}
\]
\ie $B=\{x\in A\mid\Delta_A(x)-x\widehat\otimes1\in A\widehat\otimes I\}$. Denote by $\gamma$ the continuous morphism of $k$-algebras $\theta^\natural:A'\to A$; it is a section of $\rho'$ and an isomorphism of profinite $k$-algebras from $A'$ to $B$. Moreover, $\tau=\pi^\natural$ identifies $A_0$ with a sub-bialgebra of $A$, which is precisely the affine algebra of the quotient $N\backslash G$. One then deduces from (1) and (3) that one has an isomorphism of profinite $k$-algebras
\begin{equation}\tag{*}
A'\widehat\otimes_k A_0\isomto A,\qquad a'\widehat\otimes a_0\mto\gamma(a')\tau(a_0),
\end{equation}
whose inverse is the map $a\mto(\rho'\otimes\rho)\Delta_A(a)$.
Finally, identify $A'$ with its image in $A$ under $\gamma$, so that the projection $A\to A'$ is then $\gamma\rho'$. Denoting by $\Delta_{A'}$ the comultiplication of $A'$, one then deduces from (4) that $\Delta_A(A')\subset A\widehat\otimes A'$ and that the diagram below commutes:
\[
\xymatrix{A'\ar@{=}[d]\ar[r]^{\Delta_A}&A\widehat\otimes A'\ar[d]^{\gamma\rho'\otimes\mathrm{id}}\\
A'\ar[r]_{\Delta_{A'}}&A'\widehat\otimes A'}
\]
(one therefore also has $\gamma\rho'\circ c_A=c_{A'}$, where $c_A$ (resp. $c_{A'}$) is the antipode of $A$ (resp. $A'$)). Moreover, denoting by $m_A$ multiplication on $A$, one deduces from (2) that, for every $a\in A$,
\[
\gamma\rho'(a)=(m_A\circ(\mathrm{id}\otimes\tau\rho c_A)\circ\Delta_A)(a).
\]

\numpara{2.4.B}\footnotemark[92]For simplicity, suppose that $k$ is artinian. Then the foregoing can be expressed more simply in terms of the covariant bialgebras of $G,Q,N$. Indeed, since $G=N\times Q$ as formal $k$-varieties, $H(G)=H(N)\otimes_k H(Q)$ as $k$-coalgebras. Moreover, since multiplication on $G$ is given by
\[
(x,q)\cdot(x',q')=(x\alpha(q,x'),qq'),\qquad\text{where }\alpha(q,x')=qxq^{-1},
\]
% typo?: x rather than x' is printed on the right; retained.
multiplication on $H(G)$ is given as follows: for every $x\in H(N)$ and $q\in H(Q)$,
\[
(x\otimes q)\cdot(x'\otimes q')=x\phi(q,x')\otimes qq',
\]
where $\phi:H(Q)\otimes_k H(N)$ is the morphism of $k$-coalgebras induced by $\alpha$.
% typo?: source omits the target of phi and coproduct summation in the preceding multiplication formula; retained.
Since $\alpha$ is the composite morphism below (where $\delta_G$ (resp. $m_G$) is the diagonal morphism (resp. multiplication) of $G$, $c_Q$ the inversion morphism of $Q$, and $v(q\otimes q'\otimes x)=q\otimes x\otimes q'$),
\[
Q\times N\xrightarrow{v\circ(\delta_G\times\mathrm{id})}Q\times N\times Q
\xrightarrow{\mathrm{id}\times\mathrm{id}\times c_Q}Q\times N\times Q\xrightarrow{m_G}G,
\]
one obtains, still denoting by $c_Q$ the antipode of $H(Q)$, that
\begin{equation}\tag{$\star$}
\phi(q\otimes x')=\sum_i q_ix'c_Q(q'_i)\quad\text{if}\quad\Delta_{H(Q)}(q)=\sum_i q_i\otimes q'_i.
\end{equation}
In particular, if $M$ is an abstract group and $H(Q)$ is the group algebra $kM$ (\ie $Q=\Spf^*kM$ is the \emph{constant formal $k$-group} $M_k$), for every $\gamma\in M$ and $x'\in H(N)$ one has $\phi(\gamma\otimes x')=\gamma x'\gamma^{-1}$, and this defines an action of $M$ on $H(N)$ by automorphisms of Hopf algebras.

\begin{enonce}{2.4.1. Proposition}{}\label{VIIB.2.4.1}
Let $f:G\to K$ be a morphism of formal $k$-groups. If $H=\Ker(f)$ is topologically flat over $k$, the homomorphism $f':G/H\to K$ induced by $f$ is a monomorphism.
\end{enonce}
This is a consequence of the results of Exposé IV;\nde{Cf. IV, 5.2.6.} nevertheless, we give a direct proof. Let $T$ be a formal variety of finite length over $k$ and $t$ an element of $(G/H)(T)$ such that $f'\circ t$ is the identity element of $K(T)$. We must show that $t$ is the identity element of $(G/H)(T)$. Denote by $p$ the projection $G\to G/H$ and by $X$ the fibered product $T\times_{G/H}G$.
By 2.4, $p$ is surjective and topologically flat, hence so is the morphism $\mathrm{pr}_1:X\to T$, so $\mathrm{pr}_1$ is an epimorphism by Proposition 1.3.1; it therefore suffices to show that $t\circ\mathrm{pr}_1$ is the identity element of $(G/H)(X)$. Denote by $\mathrm{pr}_2$ the projection $X\to G$; one has $t\circ\mathrm{pr}_1=p\circ\mathrm{pr}_2$, whence the equality $1=f'\circ t\circ\mathrm{pr}_1=f'\circ p\circ\mathrm{pr}_2=f\circ\mathrm{pr}_2$; then the exact sequence
\[
1\to H\to G\xrightarrow f K
\]
shows that $\mathrm{pr}_2$ factors through $H$, so $p\circ\mathrm{pr}_2$ is the zero morphism. Since $p\circ\mathrm{pr}_2=t\circ\mathrm{pr}_1$, this proves the proposition.
From the proposition one deduces the following corollary. Denote by $\OO(G),\OO(K),\OO(G/H)$ the affine algebras of $G,K,G/H$; one has seen (2.4) that $p$ induces an injection from $\OO(G/H)$ to $\OO(G)$. Moreover, by Proposition 1.3, since $f':G/H\to K$ is a monomorphism, the morphism $\OO(K)\to\OO(G/H)$ is surjective, whence:
\begin{corollary}{}
Let $f:G\to K$ be a morphism of formal $k$-groups and $H=\Ker(f)$. If $H$ is topologically flat over $k$, $\OO(G/H)$ is the image of $\OO(K)$ in $\OO(G)$.
\end{corollary}

\numpara{2.4.2}Retain the preceding notation and suppose that $H$ and $G$ are topologically flat over $k$. Then $G$ is topologically flat over $k$ and over $G/H$, so, by 1.3.3, $G/H$ is topologically flat over $k$. Consequently, by 2.4, the canonical projection $q$ from $K$ to $\Coker(f')$ is topologically flat, and $f'$ is an isomorphism from $G/H$ to $\Ker(q)$. Moreover, it is clear that $\Coker(f)=\Coker(f')$. Thus, under the hypothesis that $G$ and $H=\Ker(f)$ are topologically flat over $k$, one has obtained an isomorphism between $\Ker(q)$, the \emph{image} of $f$, and $G/\Ker(f)$, the \emph{coimage} of $f$. This implies the theorem below.
\begin{theorem}{}
Let $k$ be a field. Commutative formal $k$-groups form an abelian category.
\end{theorem}
\begin{corollary}{}
Let $k$ be a field. Affine commutative group $k$-schemes form an abelian category.\nde{See also the remarks following VI A, 5.4.3.}
\end{corollary}
This follows from the theorem and the equivalence of categories 2.2.2.

\numpara{2.5}A formal $k$-group is said to be \emph{étale} if the underlying formal variety is étale (cf. 1.6); these formal groups have a very simple structure. Indeed, suppose that $k$ is local; let $\kappa$ be the residue field of $k$, $\kappa_s$ a separable closure of $\kappa$, and $\Gamma$ the topological Galois group of $\kappa_s$ over $\kappa$. A \emph{$\Gamma$-set} consists of a set $E$ and a \emph{continuous} action of $\Gamma$ on $E$ (\ie the isotropy group of every element $x\in E$ is an open subgroup of $\Gamma$).
For every formal $k$-variety $X$, put
\[
X(\kappa_s)=\varinjlim_{\ell}X(\ell),
\]
where $\ell$ runs through the finite extensions of $\kappa$ contained in $\kappa_s$.\nde{Note that if $[\kappa_s:\kappa]=\infty$, $\kappa_s$ is not a profinite $\kappa$-algebra, so the preceding notation is an abuse of notation. Moreover, we have detailed the original in what follows.} Then $\Gamma$ acts continuously on each $X(\ell)$, hence also on $X(\kappa_s)$. Moreover, let $X_\kappa=X\widehat\otimes_k\kappa$ (cf. 1.6.C); for every $\ell$ one has $X(\ell)=X_\kappa(\ell)$, whence $X(\kappa_s)=X_\kappa(\kappa_s)$.
% typo?: action of Gamma on each X(ell) is printed even for nonnormal ell; retained.
Now suppose that $X$ is étale over $k$; then $X_\kappa$ is the formal $\kappa$-variety which is the direct sum of the $\Spec\kappa(x)$, for $x\in X$, and if one denotes by $X'_\kappa$ the $\kappa$-scheme which is the direct sum of the $\Spec\kappa(x)$, one sees that $X_\kappa(\kappa_s)$ is precisely $X'_\kappa(\kappa_s)=\Hom_{(\Sch/\kappa)}(\Spec\kappa_s,X'_\kappa)$.
Denote by $\mathcal C=(\Sch_{/\kappa}^{\text{ét}})$ the full subcategory of $(\Sch/\kappa)$ consisting of étale $\kappa$-schemes. One knows that the functor $X'\mto X'(\kappa_s)$ is an equivalence from $\mathcal C$ to the category $\mathcal C'$ of $\Gamma$-sets (cf. SGA 1, V §§ 7--8 or \cite{DG70}, § I.4 6.4); it therefore induces an equivalence between the category of $\mathcal C$-groups and that of $\mathcal C'$-groups, and one sees at once that a $\mathcal C'$-group is the same thing as an abstract group $G$ equipped with a continuous action of $\Gamma$ by group automorphisms (one then says that $G$ is a \emph{$\Gamma$-group}).
In view of the equivalences $\mathbf{Vaf}_{/k}^{\text{ét}}\isomto\mathbf{Vaf}_{/\kappa}^{\text{ét}}\isomto(\Sch_{/\kappa}^{\text{ét}})$ of 1.6.E, one therefore obtains the following:
\begin{proposition}{}
Let $k$ be a local pseudocompact ring, $\kappa$ its residue field, $\kappa_s$ a separable closure of $\kappa$, and $\Gamma=\Gal(\kappa_s/\kappa)$.
\begin{enumeratei}
\item The functor $X\mto X(\kappa_s)$ is an equivalence from the category of étale formal $k$-varieties to that of $\Gamma$-sets.
\item It induces an equivalence from the category of étale formal $k$-groups to that of $\Gamma$-groups.
\end{enumeratei}
\end{proposition}
\begin{remark}{2.5.A}\label{VIIB.2.5.A}\nde{We have added this remark.}
Let $k$ be a field, $k_s$ a separable closure of $k$, $G$ an étale formal $k$-group, and $M$ the abstract group $G(k_s)$. Denote by $X$ a set of representatives of the orbits of $\Gamma=\Gal(k_s/k)$ in $M$, and for every $x\in X$ let $\Gamma_x$ be its stabilizer, which is a subgroup of $\Gamma$ of finite index, and $L_x=k_s^{\Gamma_x}$, which is an extension of $k$ of degree $[\Gamma:\Gamma_x]$ (see, for example, \cite{BAlg}, V § 10.10). Then, by the equivalence of categories above, the affine algebra $\mathscr A(G)$ of $G$ is the product of the $L_x$, equipped with the product topology, so the $L_x$ are precisely the simple quotients $\mathscr A(G)/\mathfrak m$, where $\mathfrak m$ is an open maximal ideal of $\mathscr A(G)$. Since these correspond by duality to subcoalgebras of $H(G)$, one obtains that $H(G)$ is \emph{pointed} (\ie $\dim_k(C)=1$ for every simple subcoalgebra $C$) if and only if $L_x=k$ for every $x$; in this case $\mathscr A(G)$ is the topological algebra $k^M$, hence $H(G)$ is the group algebra $kM$, and one therefore has $M=\{x\in H(G)\mid\varepsilon(x)=1\ \text{and }\Delta(x)=x\otimes x\}=G(k)$.
\end{remark}

\numpara{2.5.1}Again suppose that the pseudocompact ring $k$ is arbitrary. Since the functor $X\mto X_e$ of 1.6.F commutes with finite products, it transforms every formal group $G$ into an étale formal group $G_e$, and since the morphism $p:X\to X_e$ of \loccit is functorial in $X$, $p:G\to G_e$ in this case is a morphism of formal groups.
\nde{We have detailed the original in what follows.}Consider the kernel $\Ker(p)$; it is the fibered product of the diagram
\[
\xymatrix{&G\ar[d]^p\\\Spf(k)\ar[r]_\varepsilon&G_e}
\]
where $\varepsilon$ is the identity section of $G_e$. Since $p$ induces a bijection of the underlying sets, one deduces (cf. 1.2, N.D.E. (41)) that the set underlying $\Ker(p)$ is the image of $\Spf(k)$ under $\varepsilon$, and that for every point $s$ of $\Spf(k)$, the local algebra of $\Ker(p)$ at the point $\varepsilon(s)$ is $\OO_{G,\varepsilon(s)}\widehat\otimes_{\OO_{G_e,\varepsilon(s)}}k_s$. Moreover, at each point $\varepsilon(s)$ the residue field of $\OO_{G,\varepsilon(s)}$ is $\kappa(s)$, and hence $\OO_{G_e,\varepsilon(s)}=k_s$, whence $\OO_{\Ker(p),\varepsilon(s)}=\OO_{G,\varepsilon(s)}$. For these reasons we shall say that $\Ker(p)$ is the \emph{infinitesimal neighborhood of the origin of $G$}, and shall write $\Ker(p)=G_0$, thus obtaining an exact sequence
\[
1\to G_0\xrightarrow{\mathrm{incl.}}G\xrightarrow p G_e.
\]
